% Part: incompleteness
% Chapter: computability-incompleteness
% Section: introduction

\documentclass[../OpenLogic/include/open-logic-section]{subfiles}

\begin{document}

\olfileid{inc}{cpi}{int}
\olsection{可计算枚举性与可公理化理论}

不完全性定理与可计算性之间有着紧密的联系。这一点从以下事实即可看出：为了得到第一和第二不完全性定理的证明论版本，我们必须先发展相当一部分可计算函数理论。我们必须定义原始递归函数，证明可以用原始递归函数来算术化可证性，并证明 $\Th{Q}$ 表示所有原始递归函数和关系，特别是任意可公理化理论~$\Th{T}$ 的证明关系~$\Prf[\Th{T}]$。不完全性定理的两个假设本质上都是关于可计算性的假设。$\Th{T}$ 是可公理化的这一假设意味着 $\Th{T}$ 由一个\emph{可判定的}公理集所公理化。而 $\Th{T}$ 扩张 $\Th{Q}$ 这一假设，则是为了保证 $\Th{T}$ 能够表示所有原始递归函数和关系。

事实上，我们可以利用可计算性理论的思想来推广不完全性定理。这只需要求我们的计算模型能够处理公式和证明。可证性的算术化表明，那些原本只直接适用于数的计算模型（通过Gödel编码）也能做到这一点。其他计算模型可能不需要这么多工作，可以免去我们为证明 $\Prf[\Th{T}]$ 的原始递归性而对序列、公式和推导进行编码的全部繁琐工作。例如，一个足够一般的Turing机概念可以让我们直接将公式和证明表示为符号序列，方法是允许磁带上出现逻辑联结词和量词、变元、谓词符号等符号。\emph{某些}编码仍然必不可少，因为Turing机的字母表只允许有限多个符号，而变元、常项符号等却有无限多个。事实上，验证存在一台判定证明关系的Turing机所需的所有细节，比证明该关系（理解为Gödel数上的关系）是原始递归的复杂得多。

$\Th{T}$ 是可公理化的这一假设可以推广为：$\Th{T}$ 的定理集是可计算枚举的。如果我们的计算模型需要Gödel编码（例如部分递归函数），那么这等价于说 \[
  \Gn{\Th{T}} = \Setabs{\Gn{!A}}{\Th{T} \Proves !A}
\] 是 c.e.\@。$\Th{T}$ 是可公理化的这一假设蕴含了这一点。

\begin{prop}\ollabel{prop:ax-theory-ce} 设 $\Th{T}$ 由一个可判定的公理集~$\Gamma$ 公理化。则 $\Th{T}$ 是 c.e.\@
\end{prop}

\begin{proof}
  非形式地说，我们可以通过在我们偏好的任何证明系统中搜索 $\Th{T}$ 的语言中所有可能的证明，来可计算地枚举 $\Th{T}$ 的所有定理。可接受的证明系统具有这样的性质：可以判定一个符号和公式的排列是否是一个正确的证明，例如，检验自然演绎中某个推理规则的给定应用是否正确是可判定的。$\Gamma$ 是可判定的这一假设使我们能够检查一个正确的推导是否是\emph{从}~$\Gamma$ 出发的推导。例如，对于自然演绎推导中的每个未解除的假设，我们可以检查它是否是 $\Gamma$ 中的公理。如果一个推导通过了这些检查，我们就输出其最终公式；它就是 $\Th{T}$ 的一个定理。

  更形式化地说，我们可以如下进行。如果 $\Gn{\Gamma}$ 是原始递归的，则 \tagrefs{prfAX/{inc:art:pax:prop:prf-prim-rec},
  prfSC/{inc:art:plk:prop:prf-prim-rec},
  prfND/{inc:art:pnd:prop:prf-prim-rec}} 表明证明谓词 $\Prf[\Gamma]$ 是原始递归的。不难验证，同样的证明也表明：如果 $\Gn{\Gamma}$ 是可计算的，那么 $\Prf[\Gamma]$ 也是可计算的。令 $f(y) = \umin{x}{\Prf[\Gamma](x,y)}$。显然，$f(y)\fdefined$ 当且仅当 $y$ 是 $\Th{T}$ 的某个定理的Gödel数。因此，$\Gn{\Th{T}}$ 是一个部分可计算函数的定义域，从而由 \olref[cmp][thy][eqc]{thm:ce-equiv} 可知它是 c.e.\@
\end{proof}

或许令人惊讶的是，\olref{prop:ax-theory-ce} 的逆命题也成立：每个 c.e.\@ 理论都是可公理化的。这要用到所谓的“Craig（克雷格）技巧”来证明。

\begin{prop}\ollabel{prop:ce-ax}
如果 $\Th{T}$ 是一个 c.e.\@ 理论，则 $\Th{T}$ 可由一个可判定的（事实上，原始递归的）公理集公理化。
\end{prop}

\begin{proof}
  由于 $\Th{T}$ 是 c.e.\@，存在它的一个可计算枚举~$f$。根据 \olref[cmp][thy][eqc]{thm:ce-equiv} 的 \olref[cmp][thy][eqc]{case:ran-prim} 情形，我们可以假设该枚举是原始递归的。令 $!A_n$ 为该枚举中的第 $n$ 个元素。考虑集合
  \[
  \Gamma = \{!A_0, !A_1 \land !A_1, !A_2 \land (!A_2 \land !A_2), \dots\}
  \]
  $\Gamma$ 显然公理化了~$\Th{T}$，因为 $\Th{T}$ 的每个定理都会作为某个 $!A_n$ 出现在枚举中，并且可以从 $!A_n \land \dots \land !A_n$（即 $!A_n$ 的 $n+1$ 个副本）推导出来，而后者正是 $\Gamma$ 的第 $(n+1)$ 个元素。$\Gamma$ 也是可判定的：要检验 $!A \in \Gamma$，先判断它是否具有 $!B \land
  (!B \land \dots \land !B)\dots)$ 的形式，如果是，再计算它包含多少个~$!B$ 作为合取支。设 $n$ 为该数目。如果 $!A$ 不是由相同合取支构成的合取式，则令 $n=0$ 且 $!B \ident !A$。那么 $!A \in \Gamma$ 当且仅当 $!B \ident !A_{n}$。

  更形式化地说，我们可以如下进行。令 $f(n) = \Gn{!A_n}$ 为 $\Th{T}$ 的一个原始递归枚举。验证满足 $i(\Gn{A}) = n$ 且 $b(\Gn{A}) =
  \Gn{B}$ 的函数~$i$ 和~$b$ 是原始递归的。$!A \in \Gamma$ 当且仅当 $b(\Gn{A}) =
  f(i(\Gn{A}))$。后一个表达式定义了一个原始递归谓词。
\end{proof}

\begin{prob}
  假设 $f$ 是某个由公式的Gödel数构成的集合~$\Gamma$ 的一个原始递归枚举。验证 \olref{prop:ce-ax} 的证明中的函数 $b$ 和~$i$ 是原始递归的。
\end{prob}

\end{document}